% Gerado a partir de MONOGRAFIA_PASE.md. O texto foi transposto sem alterações.

\chapter{Introdução}\label{cap:1}

\section{Contexto}\label{sec:1.1}

O ecossistema de software, que sustenta grande parte da economia digital, atravessa uma transformação de natureza diferente das anteriores. A computação em nuvem mudou onde o software executa e como é cobrado. A IA agêntica muda o que o software faz. Modelos de linguagem de grande escala e, mais recentemente, sistemas capazes de planejar, decidir e agir com crescente independência deslocam o software do papel de ferramenta que apoia o trabalho humano para o de sistema que executa trabalho em nome das pessoas \cite{vista2026outlook}.

A engenharia de software está no centro dessa mudança. A escrita de código, historicamente manual, evolui para uma prática de orquestração de agentes que escrevem código, avaliada e dirigida por pessoas \cite{anthropic2026trends}. A adoção já é ampla: 90\% dos profissionais de tecnologia ouvidos pelo relatório DORA de 2025 usam IA no trabalho \cite{dora2025}. Os resultados, contudo, ainda não acompanham a adoção. O mesmo relatório mostra que a IA amplifica tanto as forças quanto as disfunções das organizações, que 30\% dos profissionais confiam pouco ou nada no código gerado e que a adoção continua associada a maior instabilidade na entrega de software.

A tensão se repete em outras dimensões do ecossistema. Interfaces que se adaptam a cada usuário melhoram a usabilidade percebida, mas podem desorientar e corroer a confiança quando mudam sem explicação \cite{arora2025selfevolving}. Agentes que agem sozinhos ampliam a superfície de ataque: a maioria dos agentes avaliados em um estudo recente mostrou-se vulnerável à manipulação por instruções maliciosas (LUPINACCI \textit{et al.}, 2025 \textit{apud} CHHABRA \textit{et al.}, 2026)\nocite{chhabra2026security}. No trabalho, o uso desatento da IA reduz a formação das competências necessárias para supervisioná-la \cite{shen2026skill}.

O problema que motiva este trabalho, portanto, não é a falta de capacidade técnica. É a falta de uma arquitetura que transforme essa capacidade em algo confiável. As organizações dispõem de agentes, mas não de um plano de controle que os coordene, que preserve a decisão humana nos pontos em que a ação gera consequência, que torne o custo observável e que registre cada decisão de forma auditável \cite{isg2025agentic,spark2026cio}. Sem essa estrutura, a autonomia crescente dos agentes tende a gerar fragmentação, retrabalho e risco em vez de valor.

Este trabalho responde a esse problema com a Arquitetura PASE (\textit{Personalized Adaptive Software Ecosystem}), uma arquitetura híbrida humano-IA que estrutura o ciclo de vida completo de uma funcionalidade, da originação da ideia à sua entrega e personalização em produção, sob um regime de autonomia supervisionada. Na PASE, os agentes de IA propõem, decompõem, executam e verificam, enquanto as pessoas governam as decisões críticas e respondem por elas.

\section{Objetivos}\label{sec:1.2}

O objetivo geral deste trabalho é propor e validar a Arquitetura PASE, uma arquitetura híbrida humano-IA em três camadas, de Tarefas, de Orquestração e de Execução, que automatiza o ciclo de vida do software com agentes de IA sob supervisão humana. Os requisitos da arquitetura derivam de uma análise multidimensional do ecossistema de software, conduzida por revisão da literatura publicada entre 2017 e 2026 sobre indústria, desenvolvimento, experiência do usuário, segurança e governança, e mercado de trabalho.

Duas questões de pesquisa orientam o estudo. A primeira (QP1) consiste em investigar se a evolução da IA conduz o software a um ecossistema totalmente autônomo de agentes, que substitua progressivamente os modelos tradicionais, ou a sistemas híbridos humano-IA, que demandam novas arquiteturas, competências e formas de governança. A segunda (QP2) consiste em investigar como uma arquitetura de software pode equilibrar a autonomia crescente dos agentes de IA com a necessidade permanente de supervisão, controle e alinhamento humano, e se essa arquitetura é viável quando implementada.

Para alcançar o objetivo geral e responder a essas questões, foram definidos seis objetivos específicos. Os dois primeiros sustentam a análise: (i) analisar o impacto da IA agêntica em cinco dimensões do ecossistema de software e identificar as lacunas que separam a capacidade técnica de uma adoção confiável; e (ii) derivar dessas lacunas um conjunto de requisitos para uma arquitetura de software agêntico. O terceiro sustenta a proposta: (iii) especificar a Arquitetura PASE, seus princípios de projeto, suas camadas e seus pontos de decisão humana. Os três últimos sustentam a validação: (iv) implementar uma prova de conceito da arquitetura e aplicá-la à evolução de um produto real; (v) validá-la por meio de execuções reais, roteiros de demonstração, testes automatizados e métricas de supervisão; e (vi) discutir as limitações da proposta e os caminhos para sua evolução.

\section{Metodologia}\label{sec:1.3}

O trabalho foi conduzido em três etapas, que correspondem às três partes centrais do texto: uma revisão da literatura, a proposição de uma arquitetura e a validação dessa arquitetura por meio de uma prova de conceito. A primeira etapa consistiu em uma revisão narrativa estruturada da literatura, organizada pelas cinco dimensões analisadas no Capítulo~\ref{cap:2}. Foram incluídos trabalhos que tratam de modelos de linguagem ou de IA agêntica em pelo menos uma dessas dimensões, publicados entre 2017 e 2026, em inglês ou português, com texto integral acessível e com método declarado, fosse um estudo empírico, uma revisão ou um levantamento com amostra descrita. Foram excluídos textos opinativos sem dados, materiais promocionais sem metodologia e versões duplicadas de um mesmo trabalho. A busca considerou bases de publicações científicas de computação, como ACM Digital Library, IEEE Xplore, ScienceDirect e SpringerLink, além do repositório arXiv e de publicações oficiais de organizações de pesquisa e de mercado.

O corpus resultante reúne 22 trabalhos. Desses, nove são artigos publicados em periódicos ou eventos, oito deles em veículos com revisão por pares, e cinco são versões preliminares disponíveis no arXiv. Os demais são um relatório técnico de pesquisa industrial e sete relatórios de indústria e de mercado, como o relatório DORA, produzido a partir de uma pesquisa com 4.867 profissionais. A inclusão de versões preliminares e de relatórios de indústria foi deliberada. Os fenômenos analisados, como agentes de programação em produção e arquiteturas multiagente, intensificaram-se em 2025 e 2026, em um ritmo que os periódicos ainda não acompanham. Para equilibrar atualidade e rigor, a revisão sistemática de \citeonline{hou2024llm4se}, que cobre 395 artigos publicados desde 2017, foi usada como referência para o panorama acadêmico, e os relatórios de indústria foram empregados apenas quando declaram amostra ou método. Os trabalhos relacionados do Capítulo~\ref{cap:3} foram selecionados separadamente, entre propostas publicadas em eventos e periódicos revisados por pares.

A segunda etapa foi uma pesquisa propositiva. As lacunas identificadas na revisão foram convertidas em requisitos (Seção~\ref{sec:2.6}), e os requisitos, em princípios de projeto e componentes da arquitetura (Capítulo~\ref{cap:4}). A terceira etapa foi a validação por estudo de caso: uma prova de conceito da arquitetura foi implementada e aplicada à evolução de um produto real, um painel financeiro pessoal. Seu comportamento foi observado em execuções com modelos de linguagem reais, e as evidências foram coletadas dos registros do próprio sistema e do histórico do repositório (Capítulo~\ref{cap:5}).

\section{Organização do Texto}\label{sec:1.4}

O restante deste trabalho está organizado em cinco capítulos. O Capítulo~\ref{cap:2} analisa o software na era da IA agêntica em cinco dimensões, identifica as lacunas de cada uma e as converte nos requisitos que orientam a proposta. O Capítulo~\ref{cap:3} apresenta os trabalhos relacionados e delimita a contribuição deste estudo em relação a eles. O Capítulo~\ref{cap:4} apresenta a Arquitetura PASE, seus princípios de projeto e suas três camadas. O Capítulo~\ref{cap:5} descreve a implementação da prova de conceito e sua validação em dois fluxos de ponta a ponta, com os resultados obtidos e sua discussão. Por fim, o Capítulo~\ref{cap:6} apresenta as conclusões, as contribuições, as limitações do estudo e os trabalhos futuros.
